\section{Privacy Analysis for Matrix Mechanisms}\label{sec:cptb}


In this section, we give an algorithm for computing an upper bound on the privacy guarantees of the matrix mechanism, and prove its correctness.

\subsection{Mixture of Gaussians Mechanisms}

The key tool in our privacy analysis is a \textit{mixture of Gaussians mechanism}, a generalization of the Gaussian mechanism with sampling. Here we define these mechanisms under the add adjacency, i.e. $D'$ contains an example zeroed out in $D$.

\begin{defn}\label{def:mog}
A \textbf{mixture of Gaussians (MoG) mechanism} is defined by two lists, a list of probabilities $\{p_1, p_2, \ldots, p_k\}$, with $\sum_i p_i = 1, p_i \in [0, 1]$, a list of sensitivities $\{c_1, c_2, \ldots c_k\}$ and a noise level $\sigma$. For simplicity, we will assume $c_i \geq 0$. Given $D$, the mechanism $\calM_{MoG}(\{p_1, p_2, \ldots, p_k\}, \{c_1, c_2, \ldots c_k\})$ outputs $z \sim N(0, \sigma^2)$. Given $D'$, it samples $s$ from the distribution with support $\{c_i\}_{i \in [k]}$ and associated probabilities $\{p_i\}_{i \in [k]}$, and outputs $z \sim N(s, \sigma^2)$. In other words, it is a Gaussian mechanism where the sensitivity $s$ is a random variable distributed according to $\{p_i\}_{i \in [k]}, \{c_i\}_{i \in [k]}$. 

A \textbf{vector mixture of Gaussians (VMoG) mechanism} $\calM_{VMoG}$ is the same as a MoG mechanism, except the sensitivities $c_i$ are allowed to be vectors $\bfc_i$ instead of scalars, and our output is sampled from a multivariate Gaussian $\bfz \sim N(\boldzero, \sigma^2 \cdot \mathbb{I})$ or $\bfz \sim N(\bfs, \sigma^2 \cdot \mathbb{I})$. 
\end{defn}

It will be easier for us to work with a special case of MoG mechanisms, where the probabilities and sensitivities arise from a product distribution:

\begin{defn}
A \textbf{product mixture of Gaussians (PMoG) mechanism} is defined by two lists $\{p_1,\ldots, p_k\}$ and $\{c_1,\ldots c_k\}$ and a noise level $\sigma$. The mechanism $\calM_{PMoG}(\{p_1,\ldots, p_k\}, \{c_1,\ldots, c_k\})$ is defined equivalently as $\calM_{MoG}(\{\prod_{i \in S} p_i \cdot \prod_{i \not \in S} (1 - p_i) | S \in 2^{[k]}\}, \{\sum_{i \in S} c_i | S \in 2^{[k]}\})$.
\end{defn}

We will need a few properties about MoG mechanisms. 


\subsubsection{Monotonicity of MoG Mechanisms}
The following shows the privacy guarantees of a MoG mechanism are ``monotonic'' in the sensitivity random variable $\bfc_i$:

\begin{lem}\label{lem:vectormonotonicdominance}
Let $\{p_1, p_2, \ldots p_k\}, \{\bfc_1, \bfc_2, \ldots, \bfc_k\}$ and $\{\bfc'_1, \bfc'_2, \ldots \bfc'_{k}\}$ be such that (i) each $\bfc_i$ is non-negative and (ii) $\bfc'_i$ is entry-wise greater than or equal to $\bfc_i$ for all $i$, i.e. each $\bfc'_i - \bfc_i$ is non-negative.

Then the PLD of \[\calM_{VMoG}(\{p_1, p_2, \ldots p_k\}, \{\bfc_1, \bfc_2, \ldots, \bfc_k\})\] is dominated by the PLD of \[\calM_{VMoG}(\{p_1, p_2, \ldots p_{k}\}, \{\bfc'_1, \bfc'_2, \ldots \bfc'_{k}\}).\]
\end{lem}
\begin{proof}

By \cref{lem:addvsremove}, it suffices to only consider the remove adjacency, i.e. given $D$ we sample $\bfc_i$ and then sample from $N(\bfc_i, \sigma^2 \mathbb{I})$ and given $D'$ from $N(\boldzero, \sigma^2 \mathbb{I})$. The privacy loss of outputting $\bfx$ is:

\[PL(\bfx) := \ln\left(\sum_i p_i \exp\left(\frac{2 \langle \bfc_i, \bfx \rangle - \ltwo{\bfc_i}^2}{2\sigma^2}\right)\right).\]

Let $S \subseteq \mathbb{R}^d$ be \textit{monotonic} if for any $\bfx \in S, \bfy$ such that $\bfy - \bfx$ is non-negative, $\bfy$ is also in $S$. In other words, increasing any subset of the entries of $\bfx \in S$ gives another vector in $S$. 
Since all $\bfc_i$ are non-negative, if $\bfy - \bfx$ is non-negative, then the privacy loss of outputting $\bfy$ is larger than that of outputting $\bfx$.
So for any VMoG mechanism and any $t$, the set of outputs $S_t = \{\bfx: PL(\bfx) \geq t\}$ is monotonic. By the Neyman-Pearson lemma it suffices to consider only the sets $S_t$ in the definition of $(\epsilon, \delta)$-DP, i.e. a mechanism satisfies $(\epsilon, \delta)$-DP if and only if 

\[\forall t: \Pr_{x \sim \calM(D)}[x \in S_t] \leq e^\epsilon \cdot \Pr_{x \sim \calM(D')}[x \in S_t] + \delta.\]

So, in order to show that to show the first VMoG mechanism is dominated by the second, it suffices to show the probability that $\bfx \sim N(\bfc_i, \sigma^2 \mathbb{I})$ is in any monotonic set $S$ is at most the probability that $\bfx \sim N(\bfc'_i, \sigma^2 \mathbb{I})$ is in $S$. This is immediate by a coupling of the two random variables: we let the first random variable be $\bfc_i + \bfz$ and the second random variable be $\bfc'_i + \bfz$, where the choice of $i$ and Gaussian noise $\bfz$ are the same for both random variables. For any monotonic $S$, since $\bfc'_i - \bfc_i$ is non-negative, $\bfc_i + \bfz$ is in $S$ only if $\bfc'_i + \bfz$ is in $S$, giving that the probability $\bfx \sim N(\bfc_i, \sigma^2 \mathbb{I})$ is in $S$ is at most the probability $\bfx \sim N(\bfc'_i, \sigma^2 \mathbb{I})$ is in $S$.
\end{proof}

Since the above proof holds for any $\bfc_i, \bfc_i'$ satisfying the assumptions in the lemma, it also holds if $\bfc_i'$ are fixed/non-adaptive but the entries in $\bfc_i$ are chosen adaptively (while still satisfying the assumptions in the lemma), i.e. the $j$th coordinate of $\bfc_i$ is chosen only after seeing the first $j-1$ coordinates of the output. In the scalar case, we get the following corollary:

\begin{cor}\label{lem:monotonicdominance}
Let $\{p_1, p_2, \ldots p_k\}, \{c_1, c_2, \ldots, c_k\}$ and $\{p'_1, p'_2, \ldots p'_{k'}\}, \{c'_1, c'_2, \ldots c'_{k'}\}$ be such that for all $T$, $\sum_{i : c'_i \geq T} p'_i \geq \sum_{i: c_i \geq T} p_i$. In other words, the random variable induced by $\{p_i\}_i, \{c_i\}_i$ is stochastically dominated by the random variable induced by $\{p'_i\}_i, \{c'_i\}_i$. We also assume $c_i, c_i' \geq 0$ for all $i$.

Then the PLD of \[\calM_{MoG}(\{p_1, p_2, \ldots p_k\}, \{c_1, c_2, \ldots, c_k\})\] is dominated by the PLD of \[\calM_{MoG}(\{p'_1, p'_2, \ldots p'_{k'}\}, \{c'_1, c'_2, \ldots c'_{k'}\}).\]
\end{cor}

\cref{lem:monotonicdominance} follows from \cref{lem:vectormonotonicdominance} since by allowing duplicate $c_i$ values, we can reduce to the setting where the probabilities are the same, and $c_i \leq c_i'$ for all $c_i$. For example, if $c_i$ is 0 or 1 w.p. 1/2 and $c_i'$ is 0, 1, or 2 w.p. 1/3, we can use $\{p_i\} = \{1/3, 1/6, 1/6, 1/3\}$, $\{c_i\} = \{0, 0, 1, 1\}$, and  $\{c'_i\} = \{0, 1, 1, 2\}$.

\subsubsection{Dimension Reduction for MoG Mechanisms}

We now give the following lemma, which lets us reduce the dimensions of a VMoG mechanism.

\begin{lem}\label{lem:scalardominance}
Let $\bfc_1, \bfc_2, \ldots, \bfc_k \in \mathbb{R}^{n \times p}$. Let $\bfc_1', \bfc_2', \ldots, \bfc_k' \in \mathbb{R}^{n}$ be vectors such that $\ltwo{(\bfc_i)_{j,:}} \leq \bfc'_i(j)$ for all $i, j$, i.e. the entries of $\bfc'_i$ upper bound the $\ell_2$-norms of the corresponding rows of $\bfc_i$.
Then the PLD of \[\calM_{VMoG}(\{p_1, p_2, \ldots, p_k\}, \{\bfc_1, \bfc_2, \ldots, \bfc_k\})\] is dominated by the PLD of \[\calM_{VMoG}(\{p_1, p_2, \ldots, p_k\}, \{\bfc_1', \bfc_2', \ldots, \bfc_k'\}).\]
Furthermore, this holds even if the rows of each $\bfc_i$ are adaptively chosen and $\bfc_i'$ are fixed, i.e. the $j$th row of all $\bfc_i$ is chosen by an adversary after seeing the first $j-1$ rows of the output of the VMoG mechanism, as long as the assumption $\ltwo{(\bfc_i)_{j,:}} \leq \bfc'_i(j)$ holds.
\end{lem}

We need  the following lemma, which we can apply multiple times to prove \cref{lem:scalardominance}:

\begin{lem}\label{lem:sumofmaxes}
Let $w_1, w_2, \ldots w_k > 0$ be positive scalars and let $\bfc_1, \bfc_2, \ldots \bfc_k \in \mathbb{R}^p$ be arbitrary vectors. Then for any $\epsilon$ and $\sigma > 0$:

\[\mathbb{E}_{\bfx \sim N(\boldzero, \sigma^2 \mathbb{I}_p)}\left[\max\left\{\sum_i w_i \exp(\langle \bfc_i, \bfx\rangle) - e^\epsilon, 0\right\} \right] \leq \mathbb{E}_{x \sim N(0, \sigma^2)}\left[\max\left\{\sum_i w_i \exp(\ltwo{\bfc_i} x) - e^\epsilon, 0\right\} \right].\]
\end{lem}
\begin{proof}
$\sum_i w_i \exp(\ltwo{\bfc_i} x)$ as a function of $x$ is continuous, increasing in $x$, and has range $\mathbb{R}^+$. So, there exists some $t$ such that $\sum_i w_i \exp(\ltwo{\bfc_i} t) = e^\epsilon$. For this choice of $t$, let $t_i = w_i \exp(\ltwo{\bfc_i} t)$. Then we have for all $x$:

\[\max\left\{\sum_i w_i \exp(\ltwo{\bfc_i} x) - e^\epsilon, 0\right\} = \sum_i \max\left\{w_i \exp(\ltwo{\bfc_i} x) - t_i, 0\right\}.\]

Now, by linearity of expectation and the fact that $\max\{\sum_i a_i, \sum_i b_i\} \leq \sum_i \max\{a_i, b_i\}$:
\begin{align*}
\mathbb{E}_{\bfx \sim N(\boldzero, \sigma^2 \mathbb{I}_p)}\left[\max\left\{\sum_i w_i \exp(\langle \bfc_i, \bfx\rangle) - e^\epsilon, 0\right\} \right] &\leq \mathbb{E}_{\bfx \sim N(\boldzero, \sigma^2 \mathbb{I}_p)}\left[\sum_i \max\left\{w_i \exp(\langle \bfc_i, \bfx\rangle) - t_i, 0\right\} \right]\\
& = \sum_i \mathbb{E}_{\bfx \sim N(\boldzero, \sigma^2 \mathbb{I}_p)}\left[\max\left\{w_i \exp(\langle \bfc_i, \bfx\rangle) - t_i, 0\right\} \right]\\
& = \sum_i \mathbb{E}_{x \sim N(0, \sigma^2)}\left[\max\left\{w_i \exp(\ltwo{\bfc_i} x) - t_i, 0\right\} \right]\\
& = \mathbb{E}_{x \sim N(0, \sigma^2)}\left[\sum_i \max\left\{w_i \exp(\ltwo{\bfc_i} x) - t_i, 0\right\} \right]\\
& = \mathbb{E}_{x \sim N(0, \sigma^2)}\left[\max\left\{\sum_i w_i \exp(\ltwo{\bfc_i} x) - e^\epsilon, 0\right\} \right].
\end{align*}
\end{proof}

\begin{proof}[Proof of \cref{lem:scalardominance}]
\cref{lem:vectormonotonicdominance} holds for adaptively chosen $\bfc_i$ and fixed $\bfc_i'$ (using the notation of that lemma), so by \cref{lem:vectormonotonicdominance} it suffices to prove the lemma for adaptive $\bfc_i$ and fixed $\bfc_i'$ such that $\ltwo{(\bfc_i)_{j,:}} = \bfc'_i(j)$ for all $i, j$. Further, by~\cref{lem:addvsremove}, it suffices to show the lemma under the remove adjacency. That is, $P = N(\bfc_i, \sigma^2 \mathbb{I}_{n \times p})$, $Q = N(\boldzero, \sigma^2 \mathbb{I}_{n \times p})$, $P' = N(\bfc_i', \sigma^2 \mathbb{I}_{n})$, $Q' = N(\boldzero, \sigma^2 \mathbb{I}_{n})$, and it suffices to show $H_\epsilon(P, Q) \leq H_\epsilon(P', Q')$ for all $\epsilon$.

We have:
\begin{align*}
H_\epsilon(P, Q) &= \mathbb{E}_{\bfx \sim Q}\left[\max\left\{\frac{P(x)}{Q(x)} - e^\epsilon, 0\right\}\right]\\
&= \mathbb{E}_{\bfx \sim N(\boldzero, \sigma^2 \mathbb{I}_{n \times p})}\left[\max\left\{\frac{\sum_i p_i \exp(\ltwo{\bfx - \bfc_i}^2 / 2\sigma^2)}{\exp(\ltwo{\bfx}^2 / 2 \sigma^2)} - e^\epsilon, 0\right\}\right]\\
&= \mathbb{E}_{\bfx \sim N(\boldzero, \sigma^2 \mathbb{I}_{n \times p})}\left[\max\left\{\sum_i p_i \exp\left(\frac{2 \langle \bfc_i, \bfx \rangle - \ltwo{\bfc_i}^2}{ 2\sigma^2}\right) - e^\epsilon, 0\right\}\right]\\
\end{align*}

To reflect the fact that $\bfc_i$ can be chosen adaptively, let $\bfc_{i,j}(\bfx_{1:j-1})$ denote any adversary's adaptive choice of the jth row of $\bfc_i$ after observing the first $j-1$ rows of $\bfx$. We can then write $H_\epsilon(P, Q)$ as:

\[\mathbb{E}_{\bfx_1, \bfx_2, \ldots \bfx_n \sim N(\boldzero, \sigma^2 \mathbb{I}_{p})}\left[\max\left\{\sum_i p_i \prod_{j \in [n]} \exp\left(\frac{2 \langle \bfc_{i,j}(\bfx_{1:j-1}), \bfx_j \rangle - \ltwo{\bfc_{i,j}(\bfx_{1:j-1})}^2}{2\sigma^2}\right) - e^\epsilon, 0\right\}\right] =\]
\begin{equation}\label{eq:nestedexp}
\mathbb{E}_{\bfx_1, \bfx_2, \ldots \bfx_{n-1} \sim N(\boldzero, \sigma^2 \mathbb{I}_{p})}\left[\mathbb{E}_{\bfx_n \sim N(\boldzero, \sigma^2 \mathbb{I}_{p})}\left[\max\left\{\sum_i p_i \prod_{j \in [n]} \exp\left(\frac{2 \langle \bfc_{i,j}(\bfx_{1:j-1}), \bfx_j \rangle - \ltwo{\bfc_{i,j}(\bfx_{1:j-1})}^2}{ 2\sigma^2}\right) - e^\epsilon, 0\right\}\right]\right].
\end{equation}

Note that the values of all $\bfc_{i,j}$ in \eqref{eq:nestedexp} are constants with respect to the inner expectation. So for any realization of $\bfx_1, \bfx_2, \ldots, \bfx_{n-1}$, choosing
\[w_i = p_i \exp\left(-\frac{ \ltwo{\bfc_{i,n}(\bfx_{1:n-1})}^2}{ 2\sigma^2}\right) \prod_{j \in [n-1]}\exp\left(\frac{2 \langle \bfc_{i,j}(\bfx_{1:j-1}), \bfx_j \rangle - \ltwo{\bfc_{i,j}(\bfx_{1:j-1})}^2}{ 2\sigma^2}\right)\]
and observing that by assumption $\ltwo{\bfc_{i,n}(\bfx_{1:n-1})} = \bfc'_{i}(n)$, we can apply \cref{lem:sumofmaxes} to upper bound the inner expectation in \eqref{eq:nestedexp} as:

\[\eqref{eq:nestedexp} \leq \mathbb{E}_{\bfx_1, \bfx_2, \ldots \bfx_{n-1} \sim N(\boldzero, \sigma^2 \mathbb{I}_{p}), x_n \sim N(0, \sigma^2)}\left[\max\left\{\sum_i p_i \prod_{j \in [n-1]} \exp\left(\frac{2 \langle \bfc_{i,j}(\bfx_{1:j-1}), \bfx_j \rangle- \ltwo{\bfc_{i,j}(\bfx_{1:j-1})}^2}{ 2\sigma^2}\right)\right.\right.\]
\[\left.\left. \cdot \exp\left(\frac{2 \bfc'_i(n) x_n - \bfc'_i(n)^2}{ 2\sigma^2}\right) - e^\epsilon, 0\right\}\right].\]

We can then iteratively repeat this argument for rows $n-1$, $n-2$, \ldots $1$ to get:

\[H_{\epsilon}(P, Q) \leq \mathbb{E}_{\bfx_1, \bfx_2, \ldots \bfx_{n-1} \sim N(\boldzero, \sigma^2 \mathbb{I}_{p}), x_n \sim N(0, \sigma^2)}\left[\max\left\{\sum_i p_i \prod_{j \in [n-1]} \exp\left(\frac{2 \langle \bfc_{i,j}(\bfx_{1:j-1}), \bfx_j \rangle- \ltwo{\bfc_{i,j}(\bfx_{1:j-1})}^2}{ 2\sigma^2}\right)\right.\right.\]
\[\left.\left. \cdot \exp\left(\frac{2 \bfc'_i(n) x_n - \bfc'_i(n)^2}{ 2\sigma^2}\right) - e^\epsilon, 0\right\}\right]\]
\[\leq \mathbb{E}_{\bfx_1, \bfx_2, \ldots \bfx_{n-2} \sim N(\boldzero, \sigma^2 \mathbb{I}_{p}), x_{n-1}, x_n \sim N(0, \sigma^2)}\left[\max\left\{\sum_i p_i \prod_{j \in [n-2]} \exp\left(\frac{2 \langle \bfc_{i,j}(\bfx_{1:j-1}), \bfx_j \rangle- \ltwo{\bfc_{i,j}(\bfx_{1:j-1})}^2}{ 2\sigma^2}\right)\right.\right.\]
\[\left.\left. \cdot \prod_{j \in [n] \setminus [n-2]} \exp\left(\frac{2 \bfc'_i(j) x_j - \bfc'_i(j)^2}{ 2\sigma^2}\right) - e^\epsilon, 0\right\}\right]\]
\[\leq \mathbb{E}_{\bfx_1, \bfx_2, \ldots \bfx_{n-3} \sim N(\boldzero, \sigma^2 \mathbb{I}_{p}), x_{n-2}, x_{n-1}, x_n \sim N(0, \sigma^2)}\left[\max\left\{\sum_i p_i \prod_{j \in [n-3]} \exp\left(\frac{2 \langle \bfc_{i,j}(\bfx_{1:j-1}), \bfx_j \rangle- \ltwo{\bfc_{i,j}(\bfx_{1:j-1})}^2}{ 2\sigma^2}\right)\right.\right.\]
\[\left.\left. \cdot \prod_{j \in [n] \setminus [n-3]} \exp\left(\frac{2 \bfc'_i(j) x_j - \bfc'_i(j)^2}{ 2\sigma^2}\right) - e^\epsilon, 0\right\}\right]\]
\[\ldots\]
\[\leq \mathbb{E}_{x_1, x_2, \ldots, x_n \sim N(0, \sigma^2)}\left[\max\left\{\sum_i p_i \prod_{j \in [n]} \exp\left(\frac{2 \bfc'_i(j) x_j - \bfc'_i(j)^2}{ 2\sigma^2}\right) - e^\epsilon, 0\right\}\right]\]
\[= \mathbb{E}_{\bfx \sim N(0, \sigma^2 \mathbb{I}_n)}\left[\max\left\{\sum_i p_i  \exp\left(\frac{2 \langle \bfc_i', \bfx \rangle - \ltwo{\bfc_i'}^2}{ 2\sigma^2}\right) - e^\epsilon, 0\right\}\right] = H_\epsilon(P', Q').\]
\end{proof}

As a corollary to the above ``matrix-to-vector reduction'', we have a ``vector-to-scalar reduction'' for MoG mechanisms:

\begin{cor}\label{cor:scalardominance}
The PLD of \[\calM_{VMoG}(\{p_1, p_2, \ldots, p_k\}, \{\bfc_1, \bfc_2, \ldots, \bfc_k\})\] is dominated by the PLD of \allowbreak \[\calM_{MoG}(\{p_1, p_2, \ldots, p_k\}, \{\ltwo{\bfc_1}, \ltwo{\bfc_2}, \ldots, \ltwo{\bfc_k}\}).\]
\end{cor}

\subsection{Matrix Mechanism Conditional Composition}

\begin{figure}
\begin{algorithm}[H]
\caption{\textbf{M}atrix \textbf{M}echanism \textbf{C}onditional \textbf{C}omposition algorithm, \CPTB{}$(\bfC, p, \sigma, \delta_1, \delta_2)$}
\begin{algorithmic}[1]
\State \textbf{Input:} Matrix $\bfC$, sampling probability $p$, noise standard deviation $\sigma$, probabilities $\delta_1, \delta_2$.
\State $\{\tilde{p}_{i,j}\}_{i, j \in [n]} \leftarrow$\texttt{ProbabilityTailBounds}($\bfC, p, \sigma, \delta_1)$.\newline\Comment{$\tilde{p}_{i,j}$ is a high-probability upper bound on the probability that an example participated in round $j$, conditioned on output in rounds $1$ to $i-1$.}
\For{$i \in [n]$} 
\State $PLD_i \leftarrow$ PLD of $\calM_{PMoG}(\{\tilde{p}_{i, j}\}_{j \in [n]}, \{\bfC_{i, j}\}_{j \in [n]}).$
\EndFor
\State $PLD \leftarrow$ convolution of $\{PLD_i\}_{i \in [n]}$.
\State \Return $\min\left(\{\epsilon: PLD\text{ satisfies }(\epsilon, \delta_2)\text{-DP}\}\right)$.
\end{algorithmic}
\end{algorithm}
\vspace{-1.5em}
\caption{Algorithm \CPTB{} for computing amplified privacy guarantees of the matrix mechanism. The subroutine \texttt{ProbabilityTailBounds} is given in \cref{fig:cptb-subroutine}.}
\label{fig:cptb}
\end{figure}


\begin{figure}
\begin{algorithm}[H]
\caption{\texttt{ProbabilityTailBounds}($\bfC, p, \sigma, \delta_1)$}
\begin{algorithmic}[1]
\State \textbf{Input:} Matrix $\bfC$, sampling probability $p$, noise standard deviation $\sigma$, probability $\delta_1$.
\State $\delta'$ = $\frac{\delta_1}{2 \cdot (nnz(\bfC) - \dimension)}$ \Comment{$nnz$ is the number of non-zeros.}
\State $z = \Phi^{-1}(1 - \delta')$ \Comment{Tail bound on normal distribution; here, $\Phi$ is the standard normal CDF.}
\For{$i, j \in [\dimension]$}
\If{$\bfC_{i, j} = 0$}
\State $\tilde{p}_{i, j} = 1$
\Else 
\State $s_{i,j} = $ minimum $s$ s.t. $\Pr[\sum_{j' \leq i} x_{j'} \langle \bfC_{1:i-1,j}, \bfC_{1:i-1, j'} \rangle > s] \leq \delta', x_{j'} \stackrel{i.i.d.}{\sim} Bern(p)$\newline \Comment{$s_{i,j}$ is a tail bound on the dot product of first $i-1$ entries of $\bfC \bfx$ and $\bfC_{1:i-1,j}$.}
\State $\epsilon_{i, j} = \frac{z \ltwo{\bfC_{1:i-1, j}}}{\sigma} + \frac{2s_{i,j}- \ltwo{\bfC_{1:i-1, j}}^2}{2\sigma^2}$ \newline \Comment{$\epsilon_{i,j}$ is a tail bound on the privacy loss of a participation in round $j$ after outputting first $i-1$ rounds}
\State $\tilde{p}_{i, j} = \frac{p \cdot \exp\left(\epsilon_{i, j}\right)}{p \cdot \exp\left(\epsilon_{i, j}\right) + (1-p)}$
\EndIf
\EndFor
\State \Return $\{\tilde{p}_{i,j}\}_{i, j \in [n]}$. 
\end{algorithmic}
\end{algorithm}
\caption{Algorithm for computing $\tilde{p}_{i,j}$, tail bound on conditional probability of participating in round $j$ given first $i-1$ outputs.}
\label{fig:cptb-subroutine}
\end{figure}

The high-level idea of our algorithm, \CPTB{} (short for \textit{matrix mechanism conditional composition}), for analyzing the matrix mechanism with amplification is the following: The output of each round conditioned on the previous rounds' output is a MoG mechanism. For each round, we specify a MoG mechanism that dominates this MoG mechanism with high probability. Then by Theorem~\ref{thm:conditioningtrick}, it suffices to compute the privacy loss distribution of each of the dominating MoGs, and then use composition to get our final privacy guarantee. \MMCC{} is given in \cref{fig:cptb}. We prove that \MMCC{} computes a valid DP guarantee:

\begin{thm}\label{thm:conditioning}
Let $\epsilon$ be the output of $\CPTB{}(\bfC, p, \sigma, \delta_1, \delta_2)$. The matrix mechanism with matrix $\bfC$, uniform sampling probability $p$, and noise level $\sigma$ satisfies $(\epsilon, \delta_1 + \delta_2)$-DP.
\end{thm}
We give a high-level overview of the proof. The proof proceeds in three steps. First, we show the matrix mechanism is dominated by a sequence of non-adaptively chosen \textit{scalar} MoG mechanisms, by analyzing the distribution for each round conditioned on previous rounds and applying the vector-to-scalar reduction of \cref{lem:scalardominance} and \cref{obs:bayesianfaker}. Second, we simplify these MoG mechanisms by showing that each is dominated by a PMoG mechanism with probabilities $p_{i,j}$ depending on the outputs from previous rounds. 
Third, we show that with high probability $p_{i,j} \leq \tilde{p}_{i,j}$ for all $i,j$, i.e.,  the upper bounds generated by \texttt{ProbabilityTailBounds} hold with probability. We then apply \cref{thm:conditioningtrick}. 


\begin{proof}
For simplicity in the proof we only consider remove adjacency, i.e. $D$ contains a sensitive example zeroed out in $D'$. By symmetry the proof also works for add adjacency. By quasi-convexity of approximate DP, it suffices to prove the theorem assuming the participation of all examples except the sensitive example is deterministic, i.e. we know the contribution of all examples except the sensitive example to $\bfx$, so we can assume these contributions are zero. So, let $\bfx$ be the matrix used in the matrix mechanism if we were to sample the sensitive example in each round. Then, the matrix mechanism is a VMoG mechanism with probabilities $\{p^{|S|}(1-p)^{n - |S|}\}_{S \subseteq [n]}$ and sensitivities $\{\sum_{j \in S} \bfC_{:, j} \bfx_j\}_{S \subseteq [n]}$.

Our proof proceeds in three high-level steps: 
\begin{enumerate}
    \item We show the matrix mechanism is dominated by a sequence of adaptively chosen MoG mechanisms.
    \item We show each of the adaptively chosen MoG mechanisms is further dominated by a PMoG mechanism.
    \item We show these PMoG mechanisms are with high probability dominated by the PMoG mechanisms in~\CPTB{}, and then apply \cref{thm:conditioningtrick}. 
\end{enumerate}

\mypar{Step 1 (matrix mechanism dominated by sequence of MoG mechanisms)} 
Let $f$ be the function that takes a matrix $\bfM$ and returns a vector $f(\bfM)$ where the $i$th entry of this vector is the $\ell_2$-norm of the $i$th row of $\bfM$. Using triangle inequality, for any $\bfx$ such that each row of $\bfx$ has norm at most 1, $f(\sum_{j \in S} \bfC_{:, j} \bfx_j)$ is entrywise less than or equal to $\sum_{j \in S} \bfC_{:, j}$. So by \cref{lem:scalardominance} the matrix mechanism is dominated by the VMoG mechanism with probabilities $\{p^{|S|}(1-p)^{n - |S|}\}_{S \subseteq [n]}$ and sensitivities $\{\sum_{j \in S} \bfC_{:, j}\}_{S \subseteq [n]}$\footnote{Note that since $\bfC$ is lower-triangular, so the choice of the distribution of the $i$th row of $\bfC \bfx$ by an adaptive adversary depends only on rows $1$ to $i-1$ of $\bfC \bfx + \bfz$. That is, an adversary who chooses the $j$th row of $\bfx$ after seeing the $j-1$st first rows of the matrix mechanism satisfies the adaptivity condition in \cref{lem:scalardominance}}. Note that this is exactly the (non-adaptive) matrix mechanism where each $\bfx_i = 1$ (prior to sampling), i.e. it suffices to prove the privacy guarantee holds for this choice of $\bfx$. So, for the rest of the proof we will assume the input of the matrix mechanism (prior to sampling) is the all ones vector.

Now, let $\theta_{1:i}$ denote the output of rounds 1 to $i$. By~\cref{obs:bayesianfaker}, this random variable is the same as the composition over $i$ of outputting $\theta_i$ sampled from its distribution conditioned on $\theta_{1:i-1}$.  Let $S_i$ denote the set of rounds in $[i]$ in which we sample the sensitive example. Abusing notation to let $\Pr$ denote a likelihood, the likelihood of the matrix mechanism $\calM(D)$ outputting $\theta_i$ in round $i$ conditioned on $\theta_{1:i-1}$ is:

\[\sum_{T \subseteq [i]} \Pr_{\tau \sim \Theta(D)}[S_i = T | \tau_{1:i-1} = \theta_{1:i-1}] \Pr_{\tau_i \sim N(\sum_{j \in T} \bfC_{i,j}, \sigma^2 \cdot \mathbb{I})}\left[\tau_i = \theta_i\right]\]

The likelihood of $\calM(D')$ outputting $\theta_i$ in round $i$ (conditioned on $\theta_{1:i-1}$, which doesn't affect the likelihood since since each coordinate of $\theta$ is independent when sampled from $\calM(D')$) is

\[\Pr_{\tau_i \sim N(\boldzero, \sigma^2 \cdot \mathbb{I})}\left[\tau_i = \theta_i\right].\]

In other words, the distribution of $\theta_i$ conditioned on $\theta_{1:i-1}$ under $\calM(D), \calM(D')$ is exactly the same as the pairs of output distributions given by the MoG mechanism.

\[\calM_{MoG}\left(\left\{\Pr_{\tau \sim \Theta(D)}[S_i = T | \tau_{1:i-1} = \theta_{1:i-1}]\right\}_{T \subseteq [i]}, \{\sum_{j \in T} \bfC_{i, j} \}_{T \subseteq [i]}\right).\]

So the matrix mechanism with $\bfx$ being all ones is the same as the sequence of (adaptively chosen) MoG mechanisms given by

\[\left\{\calM_{MoG}\left(\left\{\Pr_{\tau \sim \Theta(D)}[S_i = T | \tau_{1:i-1} = \theta_{1:i-1}]\right\}_{T \subseteq [i]}, \{\sum_{j \in T} \bfC_{i, j} \}_{T \subseteq [i]}\right)\right\}_{i \in [n]}.\]

\mypar{Step 2 (each MoG is dominated by a PMoG)} To achieve step 2, we use the following lemma:

\begin{lem}\label{pmog-dominance}
Let

\[p_{i,j} = \frac{p \exp\left(\frac{2 \langle \theta_{1:i-1}, \bfC_{1:i-1, j}\rangle - \ltwo{\bfC_{1:i-1, j}}^2}{2\sigma^2}\right)}{p \exp\left(\frac{2 \langle \theta_{1:i-1}, \bfC_{1:i-1, j}\rangle - \ltwo{\bfC_{1:i-1, j}}^2}{2\sigma^2}\right) + 1 - p}.\]

The random variable induced by probabilities $\left\{\prod_{j \in T} p_{i,j} \prod_{j \in [i] \setminus T} (1 - p_{i,j})\right\}_{T \subseteq [i]}$ and support $\{\sum_{j \in T} \bfC_{i, j}\}_{T \subseteq [i]}$ stochastically dominates the random variable induced by probabilities $\{\Pr_{\tau \sim \Theta(D)}[S_i = T | \tau_{1:i-1} = \theta_{1:i-1}]\}_{T \subseteq [i]}$ and the same support.
\end{lem}

Proving \cref{pmog-dominance} completes the step as with this lemma and \cref{lem:monotonicdominance}, the PLD of 

\[\calM_{MoG}\left(\left\{\Pr_{\tau \sim \Theta(D)}[S_i = T | \tau_{1:i-1} = \theta_{1:i-1}]\right\}_{T \subseteq [i]}, \{\sum_{j \in T} \bfC_{i, j} \}_{T \subseteq [i]}\right).\]

is dominated by the PLD of

\[\calM_{PMoG}\left(\{p_{i,j}\}_{j \in [n]}, \{\bfC_{i,j}\}_{j \in [n]}\right).\]
\begin{proof}[Proof of \cref{pmog-dominance}]

Sampling $T$ according to probabilities $\{\Pr_{\tau \sim \Theta(D)}[S_i = T | \tau_{1:i-1} = \theta_{1:i-1}]\}_{T \subseteq [i]}$ is equivalent to the following process: We start with $T = \emptyset$, and for each $j \in [i]$, add it to $T$ with probability $\Pr[T \cup \{j\} \subseteq S_i | T \subseteq S_i, \tau_{1:i-1} = \theta_{1:i-1}]$. Similarly, sampling $T$ according to $\left\{\prod_{j \in T} p_{i,j} \prod_{j \in [i] \setminus T} (1 - p_{i,j})\right\}_{T \subseteq [i]}$ is equivalent to the same process, except we add $j$ with probability $p_{i,j}$. If we show that $\Pr[T \cup \{j\} \subseteq S_i | T \subseteq S_i, \tau_{1:i-1} = \theta_{1:i-1}] \leq p_{i,j}$ for all $T, j$, then we can couple these sampling processes such that with probability 1, $\sum_{j \in T} \bfC_{i,j}$ is at least as large for the second process as for the first, which implies the lemma. The posterior distribution of $S_i$ satisfies:

\[\Pr_{\tau \sim \Theta(D)}[S_i = T | \tau_{1:i-1} = \theta_{1:i-1}] \propto \Pr_{\tau \sim \Theta(D)}[S_i = T] \cdot \Pr_{\tau \sim \Theta(D)}[\tau_{1:i-1} = \theta_{1:i-1} | S_i = T]\]
\[ \propto p^{|T|}(1-p)^{i - |T|} \cdot \exp\left(\frac{2 \langle \theta_{1:i-1}, \sum_{j \in T}  \bfC_{1:i-1,j}\rangle - \ltwo{\sum_{j \in T}  \bfC_{1:i-1, j}}^2}{2 \sigma^2}\right).\]

Hence:

\[\Pr[T \cup \{j\} \subseteq S_i | T \subseteq S_i, \tau_{1:i-1} = \theta_{1:i-1}] =\]
\[\frac{\sum_{T' \supseteq T \cup \{j\}}p^{|T'|}(1-p)^{i - |T'|} \cdot \exp\left(\frac{2 \langle \theta_{1:i-1}, \sum_{j \in T'}  \bfC_{1:i-1,j}\rangle - \ltwo{\sum_{j' \in T'}  \bfC_{1:i-1, j'}}^2}{2 \sigma^2}\right)}{\sum_{T' \supseteq T}p^{|T'|}(1-p)^{i - |T'|} \cdot \exp\left(\frac{2 \langle \theta_{1:i-1}, \sum_{j \in T'}  \bfC_{1:i-1,j}\rangle - \ltwo{\sum_{j' \in T'}  \bfC_{1:i-1, j'}}^2}{2 \sigma^2}\right)}.\]

Fix some $T' \supseteq T \cup \{j\}$. Consider the term in the numerator sum corresponding to $T'$, and the two terms in the denominator sum corresponding to $T'$ and $T' \setminus \{j\}$. The ratio of the numerator term to the sum of the two denominator terms is:

\[\frac{p \cdot \exp\left(\frac{2 \langle \theta_{1:i-1}, \bfC_{1:i-1,j}\rangle-\ltwo{\sum_{j' \in T'}  \bfC_{1:i-1, j'}}^2}{2 \sigma^2}\right)}{p \cdot \exp\left(\frac{2 \langle \theta_{1:i-1}, \bfC_{1:i-1,j}\rangle-\ltwo{\sum_{j' \in T'}  \bfC_{1:i-1, j'}}^2}{2 \sigma^2}\right) + (1 - p) \cdot \exp\left(\frac{-\ltwo{\sum_{j' \in T' \setminus \{j\}}  \bfC_{1:i-1, j'}}^2}{2 \sigma^2}\right)}.\]

Since entries of $\bfC$ are non-negative, we have $\ltwo{\sum_{j' \in T'}  \bfC_{1:i-1, j'}}^2 \geq \ltwo{\sum_{j' \in T' \setminus j}  \bfC_{1:i-1, j'}}^2 + \ltwo{\bfC_{1:i-1, j'}}^2$, hence this ratio and thus $\Pr[T \cup \{j\} \subseteq S_i | T \subseteq S_i, \tau_{1:i-1} = \theta_{1:i-1}]$ are at most $p_{i,j}$, which proves the lemma.
\end{proof}

\mypar{Step 3 (replacing $p_{i,j}$ with $\tilde{p}_{i,j}$ via conditional composition)}
By \cref{thm:conditioningtrick} and \cref{lem:monotonicdominance}, it now suffices to show that w.p. $1-\delta_1$, $p_{i, j} \leq \tilde{p}_{i, j}$ for all $i, j$ simultaneously. The bound trivially holds for entries where $\bfC_{i,j} = 0$, so we only need the bound to hold for all $nnz(\bfC)$ pairs $i,j$ such that $\bfC_{i,j} > 0$. Furthermore, if $\bfC_{i,j}$ is the first non-zero entry of column $j$, then $\bfC_{1:i-1,j}$ is the all zero-vector, so we get $p_{i,j} = \tilde{p}_{i,j} = p$. 

So, there are only $nnz(\bfC) - \dimension$ ``non-trivial'' pairs we need to prove the tail bound for; by a union bound, we can show each of these bounds individually holds w.p. $\frac{\delta_1}{nnz(\bfC) - \dimension}$. By definition of $p_{i, j}, \tilde{p}_{i, j}$, this is equivalent to showing $\langle \theta_{1:i-1}, \bfC_{1:i-1, j} \rangle \leq z \ltwo{\bfC_{1:i-1,j}} \sigma + s_{i,j}$ for each of these $i, j$ pairs. We have:

\[\langle \theta_{1:i-1}, \bfC_{1:i-1, j} \rangle = \sum_{j' \in S_i} \langle \bfC_{1:i-1, j'}, \bfC_{1:i-1, j} \rangle + \langle \bfz_{1:i-1}, \bfC_{1:i-1, j}\rangle.\]

The first term is tail bounded by $s_{i,j}$ with probability $1-\frac{\delta_1}{2(nnz(\bfC) - \dimension)}$ by definition, the second term is drawn from $N(0, \ltwo{\bfC_{1:i-1, j}}^2 \sigma^2)$ and thus tail bounded by $z \ltwo{\bfC_{1:i-1, j}} \sigma$ with the same probability by definition. A union bound over these two events gives the desired tail bound on $\langle \theta_{1:i-1}, \bfC_{1:i-1, j} \rangle$.
\end{proof}

\mypar{Tightness} To get a sense for how tight \MMCC{} is, if in \MMCC{} we instead set $\tilde{p}_{i,j} = p$ for all $i,j$, this is equivalent to analyzing the matrix mechanism as if each row were independent. Since the rows are actually correlated, we expect this analysis to give a lower bound on the true value of $\epsilon$. So we can use $\max_{i,j}\tilde{p}_{i,j} / p$ as roughly an upper bound on the ratio of the $\epsilon$ reported by \CPTB{} and the true $\epsilon$ value. In particular, as $\sigma \rightarrow \infty$, for $\tilde{p}_{i,j}$ computed by \texttt{ProbabilityTailBounds} this ratio approaches 1, i.e. \CPTB{} gives tight $\epsilon$ guarantees in the limit as $\sigma \rightarrow \infty$.

\mypar{Sampling scheme of \cite{BandedMF}} The techniques used in \MMCC{} are complementary to those in \cite{BandedMF}: In \cref{sec:mmcc-extension}, we give a generalization of \MMCC{} that analyzes the matrix mechanism under their ``$b$-min-sep sampling.'' For $b = 1$, this is the same as i.i.d. sampling every round so this generalization retrieves \MMCC{}. For $b$-banded matrices this generalization retrieves exactly the DP-SGD-like analysis of \cite{BandedMF}. In other words,~\emph{this generalization subsumes all existing amplification results for matrix mechanisms.}

\mypar{Benefits of i.i.d. sampling} \MMCC{} is the first analysis that allows us to benefit from both correlated noise and privacy amplification via i.i.d. (i.e., maximally random) sampling. In \cref{sec:mse} we demonstrate that the combination of benefits allows us to get better $\ell_2^2$-error for computing all prefix sums than independent-noise mechanisms, for much smaller $\epsilon$ than prior work. 